% =============================================================================
% RÉSULTATS - Chapitre 3
% =============================================================================
% Ce chapitre présente les résultats sous forme de tableaux et figures.

\chapter{Résultats}
\label{chap:resultats}

\section{Analyse des Performances}

Les résultats sont présentés sous forme de tableaux et de figures.

\subsection{Tableaux Élégants}

Nous utilisons \texttt{booktabs} pour des tableaux professionnels (pas de lignes verticales).

\begin{table}[hbt!]
    \centering
    \caption{Comparaison des performances}
    \label{tab:perf}
    \begin{tabular}{l c c r}
        \toprule
        \textbf{Modèle} & \textbf{Précision} & \textbf{Rappel} & \textbf{Temps (s)} \\
        \midrule
        Modèle A & 0.85 & 0.78 & 12.5 \\
        Modèle B & 0.92 & 0.88 & 15.2 \\
        Modèle C & 0.89 & 0.82 & 10.1 \\
        \bottomrule
    \end{tabular}
\end{table}

\section{Illustrations}

L'utilisation de \texttt{subcaption} permet de grouper des figures.

\begin{figure}[H]
    \centering
    \begin{subfigure}[b]{0.45\textwidth}
        \centering
        \includegraphics[width=\textwidth]{img/logos/Logo_LICIT-ECO7_vertical.png}
        \caption{Logo Vertical}
        \label{fig:logo_vert}
    \end{subfigure}
    \hfill
    \begin{subfigure}[b]{0.45\textwidth}
        \centering
        \includegraphics[width=\textwidth]{img/logos/Logo_Université_Gustave_Eiffel_short.png}
        \caption{Logo Court}
        \label{fig:logo_short}
    \end{subfigure}
    \caption{Comparaison des identités visuelles}
    \label{fig:logos}
\end{figure}

Comme illustré dans la \cref{fig:logos}, nous avons deux variantes.
